\documentclass[11pt]{article}
% Shared preamble for the PNPInverse onboarding docs.
% Each document does % Shared preamble for the PNPInverse onboarding docs.
% Each document does % Shared preamble for the PNPInverse onboarding docs.
% Each document does \input{preamble} after \documentclass.

\usepackage[utf8]{inputenc}
\usepackage[T1]{fontenc}
\usepackage{lmodern}
\usepackage[margin=1in]{geometry}
\usepackage{amsmath, amssymb, amsthm}
\usepackage{mathtools}
\usepackage{empheq}
\usepackage{siunitx}
\usepackage{booktabs}
\usepackage{array}
\usepackage{enumitem}
\usepackage{xcolor}
\usepackage{framed}
\usepackage{microtype}
\usepackage{listings}
\usepackage{tikz}
\usetikzlibrary{arrows.meta, positioning, decorations.pathreplacing, calligraphy}
\usepackage[hidelinks]{hyperref}

% --- branding -------------------------------------------------------------
\definecolor{pnpblue}{RGB}{30,70,120}
\definecolor{pnpgray}{RGB}{90,90,90}
\definecolor{calloutbg}{RGB}{245,245,235}
\definecolor{shadecolor}{RGB}{245,245,235}

% --- callouts (uses framed.sty's shaded environment) ----------------------
\newenvironment{callout}[1]
  {\begin{shaded}\noindent\textbf{\color{pnpblue}#1}\par\smallskip\ignorespaces}
  {\end{shaded}}

\newcommand{\concept}[1]{\textbf{\color{pnpblue}#1}}
\newcommand{\warn}[1]{\textbf{\color{red!70!black}#1}}

% --- code listings --------------------------------------------------------
\lstset{
  basicstyle=\ttfamily\small,
  breaklines=true,
  columns=fullflexible,
  keepspaces=true,
  showstringspaces=false,
  frame=single,
  framesep=4pt,
  backgroundcolor=\color{calloutbg},
  rulecolor=\color{pnpgray!40},
}

% --- math shortcuts -------------------------------------------------------
\newcommand{\R}{\mathbb{R}}
\newcommand{\diff}{\mathrm{d}}
\newcommand{\eqd}{\overset{!}{=}}
\newcommand{\grad}{\nabla}
\newcommand{\divg}{\nabla\!\cdot}
\newcommand{\eqdef}{\;\coloneqq\;}
\newcommand{\evalat}[2]{\left.#1\right\rvert_{#2}}
\newcommand{\set}[1]{\{#1\}}

% common chemistry / physics symbols
\newcommand{\Hyd}{\mathrm{H}}
\newcommand{\Oxy}{\mathrm{O}}
\newcommand{\HtwoO}{\mathrm{H_2O}}
\newcommand{\HtwoOtwo}{\mathrm{H_2O_2}}
\newcommand{\Otwo}{\mathrm{O_2}}
\newcommand{\Hp}{\mathrm{H^+}}
\newcommand{\OHm}{\mathrm{OH^-}}
\newcommand{\Kp}{\mathrm{K^+}}
\newcommand{\Csp}{\mathrm{Cs^+}}
\newcommand{\SOfour}{\mathrm{SO_4^{2-}}}
\newcommand{\Vrhe}{V_{\mathrm{RHE}}}
\newcommand{\eo}{E^{\circ}}
\newcommand{\Eeq}{E_{\mathrm{eq}}}
\newcommand{\kzero}{k_0}

\newcommand{\onboardingfooter}{%
  \vfill\noindent\rule{\linewidth}{0.4pt}\par
  {\small\color{pnpgray} PNPInverse intern onboarding,
  compiled \today. Source under \texttt{docs/onboarding/}.}\par}
 after \documentclass.

\usepackage[utf8]{inputenc}
\usepackage[T1]{fontenc}
\usepackage{lmodern}
\usepackage[margin=1in]{geometry}
\usepackage{amsmath, amssymb, amsthm}
\usepackage{mathtools}
\usepackage{empheq}
\usepackage{siunitx}
\usepackage{booktabs}
\usepackage{array}
\usepackage{enumitem}
\usepackage{xcolor}
\usepackage{framed}
\usepackage{microtype}
\usepackage{listings}
\usepackage{tikz}
\usetikzlibrary{arrows.meta, positioning, decorations.pathreplacing, calligraphy}
\usepackage[hidelinks]{hyperref}

% --- branding -------------------------------------------------------------
\definecolor{pnpblue}{RGB}{30,70,120}
\definecolor{pnpgray}{RGB}{90,90,90}
\definecolor{calloutbg}{RGB}{245,245,235}
\definecolor{shadecolor}{RGB}{245,245,235}

% --- callouts (uses framed.sty's shaded environment) ----------------------
\newenvironment{callout}[1]
  {\begin{shaded}\noindent\textbf{\color{pnpblue}#1}\par\smallskip\ignorespaces}
  {\end{shaded}}

\newcommand{\concept}[1]{\textbf{\color{pnpblue}#1}}
\newcommand{\warn}[1]{\textbf{\color{red!70!black}#1}}

% --- code listings --------------------------------------------------------
\lstset{
  basicstyle=\ttfamily\small,
  breaklines=true,
  columns=fullflexible,
  keepspaces=true,
  showstringspaces=false,
  frame=single,
  framesep=4pt,
  backgroundcolor=\color{calloutbg},
  rulecolor=\color{pnpgray!40},
}

% --- math shortcuts -------------------------------------------------------
\newcommand{\R}{\mathbb{R}}
\newcommand{\diff}{\mathrm{d}}
\newcommand{\eqd}{\overset{!}{=}}
\newcommand{\grad}{\nabla}
\newcommand{\divg}{\nabla\!\cdot}
\newcommand{\eqdef}{\;\coloneqq\;}
\newcommand{\evalat}[2]{\left.#1\right\rvert_{#2}}
\newcommand{\set}[1]{\{#1\}}

% common chemistry / physics symbols
\newcommand{\Hyd}{\mathrm{H}}
\newcommand{\Oxy}{\mathrm{O}}
\newcommand{\HtwoO}{\mathrm{H_2O}}
\newcommand{\HtwoOtwo}{\mathrm{H_2O_2}}
\newcommand{\Otwo}{\mathrm{O_2}}
\newcommand{\Hp}{\mathrm{H^+}}
\newcommand{\OHm}{\mathrm{OH^-}}
\newcommand{\Kp}{\mathrm{K^+}}
\newcommand{\Csp}{\mathrm{Cs^+}}
\newcommand{\SOfour}{\mathrm{SO_4^{2-}}}
\newcommand{\Vrhe}{V_{\mathrm{RHE}}}
\newcommand{\eo}{E^{\circ}}
\newcommand{\Eeq}{E_{\mathrm{eq}}}
\newcommand{\kzero}{k_0}

\newcommand{\onboardingfooter}{%
  \vfill\noindent\rule{\linewidth}{0.4pt}\par
  {\small\color{pnpgray} PNPInverse intern onboarding,
  compiled \today. Source under \texttt{docs/onboarding/}.}\par}
 after \documentclass.

\usepackage[utf8]{inputenc}
\usepackage[T1]{fontenc}
\usepackage{lmodern}
\usepackage[margin=1in]{geometry}
\usepackage{amsmath, amssymb, amsthm}
\usepackage{mathtools}
\usepackage{empheq}
\usepackage{siunitx}
\usepackage{booktabs}
\usepackage{array}
\usepackage{enumitem}
\usepackage{xcolor}
\usepackage{framed}
\usepackage{microtype}
\usepackage{listings}
\usepackage{tikz}
\usetikzlibrary{arrows.meta, positioning, decorations.pathreplacing, calligraphy}
\usepackage[hidelinks]{hyperref}

% --- branding -------------------------------------------------------------
\definecolor{pnpblue}{RGB}{30,70,120}
\definecolor{pnpgray}{RGB}{90,90,90}
\definecolor{calloutbg}{RGB}{245,245,235}
\definecolor{shadecolor}{RGB}{245,245,235}

% --- callouts (uses framed.sty's shaded environment) ----------------------
\newenvironment{callout}[1]
  {\begin{shaded}\noindent\textbf{\color{pnpblue}#1}\par\smallskip\ignorespaces}
  {\end{shaded}}

\newcommand{\concept}[1]{\textbf{\color{pnpblue}#1}}
\newcommand{\warn}[1]{\textbf{\color{red!70!black}#1}}

% --- code listings --------------------------------------------------------
\lstset{
  basicstyle=\ttfamily\small,
  breaklines=true,
  columns=fullflexible,
  keepspaces=true,
  showstringspaces=false,
  frame=single,
  framesep=4pt,
  backgroundcolor=\color{calloutbg},
  rulecolor=\color{pnpgray!40},
}

% --- math shortcuts -------------------------------------------------------
\newcommand{\R}{\mathbb{R}}
\newcommand{\diff}{\mathrm{d}}
\newcommand{\eqd}{\overset{!}{=}}
\newcommand{\grad}{\nabla}
\newcommand{\divg}{\nabla\!\cdot}
\newcommand{\eqdef}{\;\coloneqq\;}
\newcommand{\evalat}[2]{\left.#1\right\rvert_{#2}}
\newcommand{\set}[1]{\{#1\}}

% common chemistry / physics symbols
\newcommand{\Hyd}{\mathrm{H}}
\newcommand{\Oxy}{\mathrm{O}}
\newcommand{\HtwoO}{\mathrm{H_2O}}
\newcommand{\HtwoOtwo}{\mathrm{H_2O_2}}
\newcommand{\Otwo}{\mathrm{O_2}}
\newcommand{\Hp}{\mathrm{H^+}}
\newcommand{\OHm}{\mathrm{OH^-}}
\newcommand{\Kp}{\mathrm{K^+}}
\newcommand{\Csp}{\mathrm{Cs^+}}
\newcommand{\SOfour}{\mathrm{SO_4^{2-}}}
\newcommand{\Vrhe}{V_{\mathrm{RHE}}}
\newcommand{\eo}{E^{\circ}}
\newcommand{\Eeq}{E_{\mathrm{eq}}}
\newcommand{\kzero}{k_0}

\newcommand{\onboardingfooter}{%
  \vfill\noindent\rule{\linewidth}{0.4pt}\par
  {\small\color{pnpgray} PNPInverse intern onboarding,
  compiled \today. Source under \texttt{docs/onboarding/}.}\par}


\title{\bfseries Phase 6$\beta$ Status \& Open Questions \\
       \large where the science is right now (May 2026)}
\author{Onboarding doc 04 of 06}
\date{\today}

\begin{document}
\maketitle

\begin{callout}{What this document is for}
Maps the recent project history onto the actual code so that the
strange-sounding labels in commit messages (``v9'', ``v10a'',
``Step~9.5'', ``Phase~D'') connect to what the model is doing and
why. Closes with a list of contributions a new collaborator can
realistically pick up.
\end{callout}

\tableofcontents
\bigskip\hrule\bigskip

% =========================================================================
\section{The arc in one paragraph}

Phase 6 is the project epoch in which we moved from a
\emph{transport-only} model that could not match the experimental
local-pH excursions to a \emph{chemistry-augmented} model that
accounts for the cation's role in modulating local proton supply.
The first sub-phase (Phase 6$\alpha$) added water self-ionization
($\HtwoO \rightleftharpoons \Hp + \OHm$) and showed it cannot
explain the deck data. The second sub-phase (Phase 6$\beta$, ongoing)
adds the actual mechanism the experimental group believes in:
field-dependent hydrolysis of hydrated alkali cations at the outer
Helmholtz plane (OHP), per Singh 2016. Phase 6$\beta$ has gone
through several internal revisions (v9, v10a, v10b) and is currently
in its first \emph{fitting} stage (Step 10 / Phase D), trying to
recover the cation-specific Singh parameters from peroxide-onset
data in the Mangan deck.

% =========================================================================
\section{Phase 6$\alpha$: water self-ionization (landed)}

\paragraph{Question.} Does adding $\HtwoO \rightleftharpoons \Hp + \OHm$
as a homogeneous bulk source/sink in the proton equation lower the
predicted surface pH from $\sim 14$ (Phase 5 baseline) into the
deck-observed window of $\sim 5$--$7$?

\paragraph{Answer.} Partially, and not enough.
\begin{itemize}[leftmargin=*]
  \item Surface pH dropped from 14.14 (transport-saturated) to 10.58
        with $K_w$ active --- a 3.5-unit improvement.
  \item The 10.58 plateau is independent of the boundary-layer
        thickness $L_{\mathrm{eff}}$: making the diffusion layer
        thinner does not lower it further. Surface pH is set by the
        \emph{kinetic-transport balance}, not by transport alone.
  \item The textbook bulk water-dissociation rate is
        $\sim 1.4$ mol/m$^3$/s; the deficit at $L = 16\,\mu$m
        cathodic demands $\sim 30$ mol/m$^3$/s --- 20$\times$
        faster than water can actually supply. So the Phase 6$\alpha$
        prediction of 10.58 is itself optimistic; the true gap to
        the deck pH window is even larger.
\end{itemize}

\paragraph{Take-away.} Water self-ionization is universal aqueous
physics and we keep it as an opt-in residual (default-off; flag
\texttt{enable\_water\_ionization}). But it is \emph{not} the
mechanism that explains cation effects. Documented in
\texttt{docs/phase6/PHASE\_6A\_INVESTIGATION\_SUMMARY.md}.

% =========================================================================
\section{Phase 6$\beta$: the cation-hydrolysis program}

\subsection{The hypothesis (recap)}

The reaction (cation hydrolysis, doc 01):
\[
   \mathrm{M(H_2O)_n^+} + \HtwoO
   \;\rightleftharpoons\;
   \mathrm{M(OH)(H_2O)_{n-1}^{0}} + \Hp,
\]
with field-dependent pK$_a$ (Singh formula, doc 01):
\[
  \mathrm{pKa}(\sigma) = \mathrm{pKa}_{\mathrm{bulk}}
     + 2 A z \sigma\, r_{H\!-\!El}
       \left(1 - \frac{r_{M-O}^2}{r_{H\!-\!El}^2}\right).
\]
This makes the cation a local proton donor whose strength depends on
both its identity (bulk pK$_a$) and the local Stern-layer surface
charge $\sigma$.

\subsection{v9: the architectural skeleton (May 2026, landed)}

v9 wired this into the solver. Specifically:
\begin{itemize}[leftmargin=*]
  \item Promoted $\Kp$ from an analytic Bikerman counter-ion to a
        \emph{dynamic} NP species (constants
        \texttt{FOUR\_SPECIES\_LOGC\_DYNAMIC\_K2SO4} in
        \texttt{scripts/\_bv\_common.py}); $\SOfour$ stays analytic.
        Reason: the Boltzmann profile gives zero NP flux by
        construction, so the kinetic source from
        $c_M(0) \cdot k_{\mathrm{hyd}}$ has nothing to drain.
  \item Added a global Real-element scalar $\Gamma_{\mathrm{MOH}}$ (mol/m$^2$)
        on the electrode, evolving by the coverage-rate ODE
        (doc 02, $\dot\Gamma$ formula).
  \item Implemented Singh's pK$_a$ formula in
        \texttt{Forward/bv\_solver/cation\_hydrolysis.py} (function
        \texttt{singh\_pka\_field\_dependent}).
  \item Outer Picard for $\Gamma_{\mathrm{MOH}}$ around the inner
        Newton solve.
\end{itemize}

\paragraph{What v9 found.} The architecture works, but the v9
formula has \emph{no Langmuir cap} on $\Gamma_{\mathrm{MOH}}$.
Converged solutions at $k_{\mathrm{hyd}} \ge 10^{-3}$ pile up
$\sim 6$ monolayers of MOH$^{0}$ at the OHP --- physically absurd
since you cannot have more than one monolayer of adsorbed species
without a second layer beneath it. v9 is therefore a
\emph{numerical} achievement, not a physically defensible model.
Documented in
\texttt{docs/phase6/PHASE\_6B\_V9\_PHASES\_A\_B\_RESULTS\_2026-05-10.md}.

\subsection{v10a: the Langmuir cap (May 2026, landed)}

v10a multiplied the hydrolysis-forward rate by
$(1 - \Gamma/\Gamma_{\max})$ so $\Gamma$ saturates at one monolayer
worth of MOH coverage:
\[
  \dot\Gamma = k_{\mathrm{hyd}} (1 - \Gamma/\Gamma_{\max})\,c_M(0)
             - k_{\mathrm{des}} \Gamma.
\]
This restored physical sensibility and unlocked the next round of
calibration questions: \emph{what is $\Gamma_{\max}$? what is
$k_{\mathrm{hyd}}$? what is $k_{\mathrm{des}}$?}

\subsection{v10a$'$ and Phase A.2: pick a calibration voltage}

The current-voltage curve has many points; if you fit cation
parameters using all 24 points at once, you cannot tell which point
the residual is being driven by. v10a$'$ proposed a
\concept{kinetic-anchor voltage} $V_{\mathrm{kin}}$ where (i) the
current is well in the kinetic-controlled regime (not Levich-limited),
(ii) the Stern surface charge is large enough that cation hydrolysis
matters, and (iii) the multi-ion solver actually converges from
cold-start. The chosen value: $V_{\mathrm{kin}} = -0.10$ V vs.\ RHE.

Phase A.2 then re-derived a clean baseline at $V_{\mathrm{kin}}$
with v10a's Langmuir cap, providing a high-confidence anchor for
later sweeps. Driver:
\texttt{scripts/studies/phase6b\_v10a\_phase\_A2\_v\_kin.py}.

\subsection{v10b: literature anchoring of constants (May 2026, landed)}

v10b set the three free Phase 6$\beta$ constants to their
literature-anchored values:
\begin{center}
\begin{tabular}{lll}
\toprule
\textbf{Constant} & \textbf{Value} & \textbf{Source} \\
\midrule
$\Gamma_{\max}$ (nondim) & 0.047 & V10A chain (tightened); literature compatibility audit (4 cross-tests). \\
$k_{\mathrm{des}}$ (nondim) & 1.0 & Engineering choice; Eyring prior $[10^{-2}, 10^{2}]$ $\Leftrightarrow$ $\Delta G_{\mathrm{des}} \in [0.69, 0.94]$ eV. \\
$C_S$ (F/m$^2$) & 0.20 & Bohra--Koper--Choi consensus on CMK-3 carbon. \\
\bottomrule
\end{tabular}
\end{center}
The constants live in the top-level Firedrake-free
\texttt{calibration/v10b.py}, importable without the heavy stack.
\warn{Deprecation alias:} \texttt{SMOKE = V10A\_SMOKE} (NEVER
\texttt{SMOKE = V10B}). Documented in
\texttt{docs/phase6/v10b\_calibration\_summary.md}.

\subsection{Steps 6--9: ablation, audit, sweeps}

After v10b landed, the project enumerated a series of
\emph{step-numbered} verification and sweeping tasks that lock in
plumbing-level correctness before a fitting attempt:

\begin{description}[leftmargin=*]
  \item[Step 6: ablation flags.] Five default-off knobs in
       \texttt{Forward/bv\_solver/config.py} that let you turn
       individual physics terms on or off
       (\texttt{apply\_h\_source}, \texttt{apply\_k\_sink}, etc.).
       All five ablations passed at A.2 baseline.
  \item[Step 7: $C_S$ literature lock.]
       \texttt{docs/phase6/CMK3\_capacitance\_literature.md} is the
       citation chain that pins $C_S = 0.20$ F/m$^2$.
  \item[Step 8: missing-data ledger.] Open data items the project
       still needs from the experimental group; tracked in
       \texttt{docs/phase6/missing\_data.md}.
  \item[Step 9: $k_{\mathrm{hyd}} \times \lambda$ ramp at v10b
       parameters.] A 14$\times$10 = 140-rung sweep over hydrolysis
       rate and the activation parameter $\lambda \in [0,1]$ at
       $V_{\mathrm{kin}}$. \textbf{Result (2026-05-11):} 14/14
       converge with mass balance $\le 5\!\times\!10^{-13}$;
       baseline reproduction byte-equivalent to v10b A.2.
       Documented in
       \texttt{docs/phase6/phase6b\_step9\_B2\_summary.md}.
  \item[Step 9.5: \texttt{warm\_start\_floor} in AdaptiveLadder.]
       Required to converge the three highest-$k_{\mathrm{hyd}}$
       extension points in step 9. Opt-in floor; default-off behavior
       is byte-equivalent. Documented in the same summary.
\end{description}

\subsection{Step 10 / Phase D: derivative-free fit of $\Delta_\beta$
                                                (in flight)}

\textbf{This is where the project is right now} (May 2026).

\paragraph{Question.} Given the v10b-anchored model with all other
constants pinned, what value of the free coupling parameter
$\Delta_\beta$ (the additive shift to Singh's $\beta_M = -A z r_{H-El}
(1 - r_{M-O}^2/r_{H-El}^2)$ per cation) best reproduces the
peroxide-onset voltage in the Mangan deck for $\Kp$ at pH 4?

\paragraph{Status.} Plan was iterated through 7 critique rounds
(\texttt{docs/handoffs/CHATGPT\_HANDOFF\_37\_phase6b-step10-phase-D/});
APPROVED at round 7. Two structural corrections in early rounds:
\begin{enumerate}[leftmargin=*]
\item the original $\Delta_\beta$ bracket was off by 7 decades
      (units: $\beta$ is in pm$^{2}$, not 1/pm$^{2}$);
\item the original continuation topology was invented; switched to
      the validated A.2 topology (anchor at $\lambda=0$, V warm-walk
      at $\lambda=0$, per-V $\lambda$-ramp $0 \to 1$). Wall budget
      grew from 5 hours to 15 hours --- overnight execution scope.
\end{enumerate}

Phase D code is in active development:
\begin{itemize}[leftmargin=*]
\item \texttt{scripts/studies/phase6b\_step10\_phase\_D\_orchestrate.py}
\item \texttt{scripts/studies/phase6b\_step10\_phase\_D\_fit\_eval.py}
\item \texttt{scripts/studies/\_phase\_D\_validate\_optimized\_path.py}
\item Tests:
      \texttt{tests/test\_phase6b\_step10\_phase\_D\_\{orchestrate,fit\_eval,plumbing\}.py}.
\end{itemize}

The optimizer is \texttt{scipy.optimize.minimize\_scalar} with
\texttt{method="bounded"} (Brent's method), since adjoints through
the Picard outer loop are not yet hooked up. The bounded interval
in target-$\Delta\mathrm{pKa}$-effect space gives a one-sided search
respecting the negative-$\Delta_\beta$ sign-guard requirement.

% =========================================================================
\section{What is locked vs.\ what is open}

\subsection{Locked (don't change without conversation)}

\begin{itemize}[leftmargin=*]
  \item Production stack: 3 dynamic species + analytic Bikerman
        counter-ions, \texttt{logc\_muh}, log-rate BV with parallel
        2e/4e at Ruggiero $\eo$.
  \item $C_S = 0.20$ F/m$^2$ (Step 7 lock; CMK-3 literature
        consensus).
  \item $\Gamma_{\max} = 0.047$ (nondim) and
        $k_{\mathrm{des}} = 1.0$ (nondim) (v10b lock).
  \item Hard rules \#1--\#6 from CLAUDE.md (anchor-and-grid,
        \texttt{exponent\_clip = 100}, IC/residual saturation
        agreement, parallel 2e/4e topology, Singh 2016 mechanism,
        $C_S$ value).
  \item Acceptance bundle (the gate set Phase D will be judged
        against): primary metric $\HtwoOtwo$\% within $\pm 10$ pp of
        deck $\Kp$ pH 4; secondaries on disk current and surface pH
        envelopes. Documented in
        \texttt{docs/phase6/PHASE\_0\_ACCEPTANCE\_BUNDLE\_LOCK\_2026-05-10.md}.
\end{itemize}

\subsection{Open}

\begin{itemize}[leftmargin=*]
  \item Phase D fit completion across all four cations (the current
        scope is $\Kp$-only first; $\mathrm{Li^+}$, $\mathrm{Na^+}$,
        $\Csp$ are queued).
  \item Adjoint annotation through the cation-hydrolysis Picard
        loop (currently a derivative-free fit; to be replaced when
        the inverse pipeline restarts).
  \item Phase 6$\alpha$.1: a finite-rate water-dissociation
        refinement, since textbook $K_w$ relaxation is too slow to
        feed the surface at the rate the equilibrium model assumes.
        Queued, not started.
  \item Per-cation Tafel slope extraction from Brianna's 2019 LSV
        data (\texttt{scripts/derive/extract\_k\_plus\_tafel\_slopes.py};
        K$^+$ done, others pending).
  \item Comparison plotting between simulated and deck observables
        (currently ad-hoc; no canonical figure-generation script).
  \item Test coverage for the new Phase D helpers
        (\texttt{tests/test\_phase6b\_step10\_phase\_D\_*.py} are
        scaffolds; flesh out parameterizations).
\end{itemize}

% =========================================================================
\section{Where a new contributor can plug in}

In rough increasing order of scope:

\begin{enumerate}[leftmargin=*]
\item \textbf{Read the chronology and ask questions.} You are coming
      in mid-stream; spend the first week on docs 00--04 and the
      acceptance-bundle lock. Bring a list of things you don't
      understand to our next 1:1.

\item \textbf{Run the Phase A.2 baseline yourself.} Activate the
      Firedrake venv (doc 05), run
      \texttt{phase6b\_v10a\_phase\_A2\_v\_kin.py}, look at the JSON
      and the plot. Verify it matches the existing
      \texttt{StudyResults/} subdirectory.

\item \textbf{Add a small diagnostic.} For example: a plot of
      $\Gamma_{\mathrm{MOH}}$ vs.\ $\Vrhe$ overlaid on the BV-rate
      contributions, or a histogram of pK$_a$($\sigma$) over the
      deck V-band. The pattern is in doc 03 \S 6.

\item \textbf{Reproduce a Singh table number.} Pick a row of the
      Singh Table S1 (Li, Na, K, Rb, Cs); plug $r_M$ and $z$ into
      the bulk pK$_a$ formula (Singh formula with $\sigma=0$);
      verify our \texttt{calibration/singh2016.py} returns the same
      number to four decimal places. This is exactly the kind of
      numerical-verification test the project values most.

\item \textbf{Improve a regression test.} The Phase D scaffolding
      tests are minimal; pick one and add three or four
      property-based parameterizations (e.g.\
      ``$\Delta_\beta = 0$ should reproduce the v10b A.2 baseline
      byte-for-byte''; ``$\Delta_\beta \to \infty$ should saturate
      $\Gamma_{\mathrm{MOH}}$ at $\Gamma_{\max}$''; etc.).

\item \textbf{Independently audit a Phase D handoff round.} The
      seven-round critique on the Phase D plan is checked into
      \texttt{docs/handoffs/CHATGPT\_HANDOFF\_37\_*}. Read R1 and
      R2; can you spot the dimensional-bracket error \emph{before}
      reading R1's accepted issues? This trains you to read the
      acceptance ledgers.
\end{enumerate}

After 4--6 weeks of these we will pick something more ambitious from
the ``open'' list above.

\onboardingfooter
\end{document}
